\chapter{THERMAL ANALYSIS BACKGROUND}
\label{app:thermal_background}
Assuming energy balance
\begin{align}
    \dot{E}_{\mathrm{system}}=\dot{E}_{\mathrm{in}}-\dot{E}_{out}\qc
\end{align}
the first law of thermodynamics 
\begin{align}
    \Delta E_\mathrm{system}=Q-W\qc
\end{align}
can be used to relate heat fluxes in and out of a system to each other. Here, $Q$ is the net heat added to the system and $W$ is the work done on the system. Assuming that no work is being done (no volume expansion), and that only the internal energy (as opposed to e.g. kinetic energy) is of interest, this becomes
\begin{align}
    \Delta {E}_\mathrm{system}&=Q\\ 
    &=m_s\cdot c_{p,s} \cdot \Delta {T}_s\qc
\end{align}
where the subscript $_s$ signals system parameters, $m$ is the mass, $c_p$ the specific heat capacity and $T$ the temperature. Finally, this gives 
\begin{align}
    m_s\cdot c_{p,s} \cdot \dot{T}_s &= \dot{E}_{\mathrm{in}}-\dot{E}_{out}\\ 
    &= \dot{Q}_\mathrm{in}-\dot{Q}_\mathrm{out}
\end{align}
For simplicity, in the following the $\dot{}$ is dropped from $Q$. 

Potential contributors to the heat flux into the system $Q_\mathrm{in}$ are 
\begin{itemize}
    \item heat generated by electronics 
    \item heat generated by pumps 
    \item heat generated by pressurisation of air 
    \item heat through absorbance of solar radiation,
\end{itemize}
while contributors to heat flux out of the system $Q_\mathrm{out}$ are 
\begin{itemize}
    \item convective cooling via ambient air
    \item radiative cooling via emittance of infrared radiation.
\end{itemize}
In the following, a derivation of the individual contributions will follow. 

\paragraph{The heat generated by compressed air} must be analysed in two stages, as there are two compression stages in the system. The first is facilitated by the vacuum pumps, pressurising ambient air to 1 bar, the second by the compressor, pressurising to 3 bar. The discharge temperature is defined as 
\begin{align}
    T_\mathrm{out}&=T_\mathrm{in}\qty(\frac{p_\mathrm{out}}{p_\mathrm{in}})^{\frac{n-1}{n}}\qc
\end{align}
where $T_\mathrm{out}$ is the discharge temperature (air temperature leaving the pumps), $T_\mathrm{in}$ is the temperature of the air on the suction side, $p_\mathrm{in}$ is the air pressure on the suction side and $p_\mathrm{out}$ is the discharge pressure. $n$ is the polytropic factor, determined by the pump and gas. If the process were isentropic, the factor would be equal to the isentropic factor of air $\gamma=1.4$. For the diaphragm pump, due to negligible friction, $n<\gamma$ as the air is dumping part of its heat into the pump.
Conversely, the compressor, as a piston pump, might generate significant heat, further heating the air, such that $n>\gamma$.
Although the discharge temperature can reach up to 500°C (assuming $n=\gamma$), the actual heat capacity of that air diminishes, as the air density decreases. Thus the total heat energy carried by the air can be calculated as  
\begin{align}
    Q_\mathrm{heat} &= \dot{m}c_p\Delta T\\ 
    &= \rho I_\mathrm{pump} c_p\Delta T\qc
\end{align}
where $\rho$ is the air density, $I_{pump}$ the volumetric pump flow rate, $c_p$ the specific heat of air, and $\Delta T$ the temperature difference between suction side and discharge side. To determine the final discharge temperature, it needs to be calculated for the diaphragm pumps and based on that for the compressor.

\paragraph{The heat generated by the pumps} themselves can be estimated from their power requirements and an efficiency factor. The diaphragm pumps will probably have a higher efficiency than the piston compressor pump, since they generate less friction. The total heat is then estimated as 
\begin{align}
    Q_\mathrm{heat} &= \epsilon P_\mathrm{draw}\qc
\end{align}
where $\epsilon$ is the efficiency factor ranging between 0 and 1 and $P_\mathrm{draw}$ is the electrical power consumption of the pump. Note, that the operations schedule implies $P_\mathrm{draw}$ to vary with time.

\paragraph{The heat generated by other electric components} is similarly calculated. Given in \tref{tab:power_full_flight_duty} is the total power draw from the BEXUS system, as well as the total consumption. This is already assuming an efficiency factor, such that the difference between those values is assumed to be emitted as heat. Thus,
\begin{align}
     Q_\mathrm{heat} &= P_\mathrm{draw}- P_\mathrm{consumed}   
\end{align}

\paragraph{The heat absorbed through solar radiation} depends on the system material, surface, and dimensions, as well as on weather conditions like cloudiness and time/solar zenith angle and altitude. In general the total heat flux is comprised of three components: direct flux from the sun, secondary solar flux reflected from Earth, clouds, and atmosphere and the heat of the earth in the \gls{ir}. Thus,
\begin{align}
    q &= q_\mathrm{sun}+q_\mathrm{reflected}+q_\mathrm{IR\,earth}\\
    &=\alpha_s (I_s \cos{\theta}+ I_s a F_{se} ) + \alpha_\mathrm{IR}(q_\mathrm{earth} F_{se})\qc
\end{align}
where $\alpha_s$ is the solar absorptivity of the surface, $I_s$ is the solar intensity received, $\theta$ is the angle between surface normal and sun, $a$ is the albedo coefficient of the Earth and atmosphere, $F_{se}$ is the view factor relating the amount of surface viewing in direction of the Earth, $\alpha_\mathrm{IR}$ is the \gls{ir} absorptivity of the surface, and $q_\mathrm{earth}$ is the \gls{ir} heat flux emitted by the Earth. Most of these values are or can be assumed to be roughly constant with respect to height. As the chassis and chamber are made out of aluminum the following analysis assumes $\alpha_{IR}=\dos{\epsilon}=0.1$ and $\alpha_s=0.4$.
The most dramatic change with altitude is seen in $I_s$.\par 
The atmosphere is absorbing part of the solar radiation incident on the top atmospheric layers $G_{sc}$. This can be modeled via the Lambert-Beer law as 
\begin{align}
    I_s(h) = G_{sc}\tau_0^{\frac{p(h)}{p_0}\frac{1}{\cos{\phi}+  0.505072 * (96.07995 - \phi) ^ {-1.6364}
}}\qc
\end{align}
where $\tau_0$ is the transmittance of solar light at ground level, $p_0$ is the atmospheric pressure at ground level, $p(h)$ is the pressure at a given altitude, and $\phi$ is the zenith angle. This formula assumes zero clouds. Since the minimum expected zenith angle in Kiruna in October is 77.11°, this value will be used by the following analysis.\par 
The final heat load is determined by multiplying the heat load for individual surfaces and summing over them, i.e.
\begin{align}
    Q_\mathrm{heat}=\sum_i q_i A_i.
\end{align}

\paragraph{The heat lost via convective cooling} via ambient air depends on the ambient air density and on system geometry. In general, 
\begin{align}
    Q_\mathrm{heat,A} = hA(T_\mathrm{system}-T_\mathrm{ambient})\qc
\end{align}
where $A$ is the area of the respective surface, $h$ is the convection coefficient and $T_\mathrm{ambient}$ is the temperature of the ambient air. The convection coefficient depends on the air density and can be modelled as 
\begin{align}
    h = h_0 \qty(\frac{p}{p_0})^n\qc
\end{align}
where $h_0$ defines some reference convection coefficient on the ground and is usually $h_0=15\,\mathrm{W/m^2/K}$. 
$n$ is usually set to 0.5 for balloon applications with laminar flow (true during ascend and floating phase, if the wind speed is not extremely high). The total heat load is again defined by adding the individual heat loads of the surfaces together.

\paragraph{The heat lost via blackbody radiation} is defined via the Stefan-Boltzmann law 
\begin{align}
    Q_\mathrm{heat,A}&= A\varepsilon\sigma (T_\mathrm{system}^4-T_\mathrm{ambient}^4)\qc
\end{align}
where $\varepsilon$ is the emissivity and $\sigma$ is the Stefan-Boltzmann constant. The total heat load is again defined by adding the individual heat loads of the surfaces together.

\paragraph{The lumped system approach} is used to model subsystems with thermally distinct behaviours. Here, this means that the system is split thermally into two parts: 
\begin{itemize}
    \item the box/chassis containing the electronics and components, 
    \item the pressure chamber within,
    \item and potentially isolation if considered (see cold and hot case in \sref{sec: thermal cold case} and \sref{sec: thermal hot case}, respectively).
\end{itemize}
The division of the system into the lumps chassis and chamber is done as the pressure chamber is to be kept at a different temperature than the rest of the system. First, the temperature of the compartment/chassis containing all electronics and components is calculated based on a constant temperature of the pressure chamber. Then, the heat flux necessary to keep the pressure chamber at the desired temperature is determined. This means that the heat flux between chamber and chassis is added as a heat load to both. 
This transfer flux is defined as 
\begin{align}
    Q_\mathrm{transfer} = F_g (T_\mathrm{chamber}-T_\mathrm{chassis}) + \epsilon_\mathrm{chamber}\sigma A_\mathrm{chamber}\qty(T_\mathrm{chamber}^4-T_\mathrm{chassis}^4)
\end{align}
Here, $F_g$ is a coefficient determining the efficiency of conductive and radiative heat transfer between chassis and chamber. While this coefficient depends on the exact geometry and material of conncecting bolts and surfaces, as well as on the air density inside the chassis, it is approximated at $F_g\approx 1.5$ W/K. The second term is responsible for modeling the radiative heat transfer between chamber and chassis. As the chassis has a much greater surface area than the chamber, it is modelled as a heat sink. This means that the efficiency of the radiative transfer is solely determined by chamber emissivity $\epsilon$ and surface area $A$.

\paragraph{Insulation} is modelled as a white styrofoam box of 5 cm thickness, completely encapsulating the chassis. The chassis is no longer subject to solar radiation, air convection, and can no longer loose heat via blackbody radiation. Thus, $Q_\mathrm{rad\,in}=Q_\mathrm{rad\,out}=Q_\mathrm{conv}=0$. Instead, the chassis is interacting with the styrofoam box via 
\begin{align}
    Q_\mathrm{insulation}=\kappa_\mathrm{styro}\frac{A_\mathrm{styro}}{d_\mathrm{styro}}\qty(T_\mathrm{styro}-T_\mathrm{chassis})\qc
\end{align} 
where $\kappa_\mathrm{styro}$ is the thermal conductivity and set to be 0.033 $\mathrm{W\,m^{-1}\,K^{-1}}$, $A_\mathrm{styro}$ is the insulation surface area, $d_\mathrm{styro}$ its thickness and $T_\mathrm{styro}$ its outer temperature. The outer temperature in turn is defined via 
\begin{align}
    Q_\mathrm{rad\,in} = Q_\mathrm{rad\,out}+Q_\mathrm{conv} + Q_\mathrm{insulation}.
\end{align}
This means that a steady-state temperature of the styrofoam is assumed. This holds true if the styrofoam's thermal capacity is small compared to that of the chassis. Since styrofoam is lightweight and the chassis is made out of aluminium this is justified.
The \gls{ir} emissivity is set to $\epsilon=\alpha_\mathrm{IR}=0.6$ and the solar absorptivity to $\alpha_s=0.15$.\par 

All in all this means that for the steady state assumptions used in the cold and hot cases, the following two equations need to be solved:
\begin{align}
    Q_\mathrm{ele}+Q_\mathrm{pump}+Q_\mathrm{transfer} + Q_\mathrm{rad\,in} + Q_\mathrm{insulation} &= Q_\mathrm{conv} + Q_\mathrm{rad\,out} \label{eq:heat_balance_compartment}\\ 
    Q_\mathrm{thermal\,system}+Q_\mathrm{compressed\,air} &=  Q_\mathrm{transfer} \label{eq:heat_balance_chamber}
\end{align} 
\eqref{eq:heat_balance_compartment} describes the heat balance of the bigger chassis and \eqref{eq:heat_balance_chamber} describes the heat balance of the pressure chamber within. Note, that depending on the assumptions in the cold and hot cases, some of the terms might be 0.